%\addcontentsline{toc}{chapter}{Socio-economical context}
\section{Socio-economical context}
\label{context}
Europe's population is ageing.
This phenomena occurs because of the augment in the life expectancy and the declining of birth rates. 
Recent studies provided by the World Health Organization (WHO) showed that the life expectancy in Spain is the highest in Europe, specially in women \cite{who_live_expectancy}.
Spanish women live an average of  85.1 years in 2012, and the men around 79 years. 
In fact, the Spanish women life expectancy is only behind the Japanese women, that is, the second highest level in the world. 
\\

In the World Population Ageing Report of 2013 it is stated that "by 2050, Bosnia and
Herzegovina, Germany, Malta, Portugal, Serbia and Spain are projected to attain median ages of
50 years or more"\cite{world_population_ageing}. 
This means that soon enough there are going to be more elder people than young people. 
How are we going to be able to support and aid that major sector of the population?
The solution might be to incorporate robots to help in this task \cite{Broekens2009}. 
\\

Assistive robots can help elder people to interact with their environment. 
They can cooperate in the health control of the user testing the blood pressure or recognizing falls. 
Also, social robots may maintain a conversation and interact with people, for example, reminding them to take a specific medication. 
Japan is the country with the highest population ageing of the world. 
This country is the world leader in the assistive robotics field. 
Recently the he International Organization for Standardization (ISO) has adopted Japanese standards for assistive robotics technology \cite{japan1}\cite{japan2}.
The Japanese government is studying to introduce a system of assistive technologies and hence the funding in this field is high. 
As an example of the technology being developed in this country, the Pepper assistive robot was recently anounced by SoftBank Corp\cite{cheap_humanoid}.
It is a robot that costs less than two thousand USD and that is capable of maintaining a conversation with the user. 
However, Pepper is not capable nowadays of performing assistive tasks. 
It has a humanoid form with arms, but they are only used to express and empathize the information that it is giving. 
Is, hence, the first step towards more environment-aware and interactive low-cost robotics.
% Is humanoid and has arms, but at least until now they are only used to express and empathize its words.  
\\

In Spain the assistive field is mainly being evolved in universities and research groups. 
Projects such as the RobAlz between the Fundación Alzheimer España (FAE) and the Carlos III University \cite{robalz} are developed to improve the interaction between robots and, in this particular case, Alzheimer patients.  
Another example of the Spanish research in the field is the work performed by the Ave María Foundation \cite{assistive_spain}.
They have introduced various robots in the therapies of their patients such as the NAO or the REEM robots. 
NAO is a small robot that can recognize faces, sounds and is able to walk. 
It is used to instruct therapy seasons. 
The REEM robot, on the other hand is used to transport the necessary items such as clothes or food around the place, allowing the employees to focus on assistive tasks. 
Nevertheless, the research in Spain is severely restricted due to the dramatic reduction that has been inflicted in its budget over the past five years. 
Recently, the Spanish government presented the 2014's Research and Innovation budget.
The numbers show that in 2014 the budget is of 6.148 millions of euros, almost a third less than the budgets of 2008, 2009 and 2010 \cite{budget_i_mas_d}.
\\

The computer vision field has numerous applications in very different industries. 
Those range from assistive robotics and medicine to industrial and defense and security programs or even autonomous cars.
It is a very active field in which many big companies such Bosch, Siemens or Sony are investing. 
In Spain the computer vision research is mainly performed in the universities in projects such as \cite{robalz}. 
It is used as a tool for the construction of bigger robotic systems. 
Nevertheless, it is affected by the low budget available for I+D in this country as well, a fact that reduces the number of research projects being developed. 
% \\

% The effect of this reduction on the research budget impulses the development of low-cost projects. 
% Also it increases the importance of using Open Source code since it is free resource. 
% In this context, this thesis is


% In the past seven years the world has experienced the effects of the recession that appeared in the late 2000s decade. It still affects the economy nowadays and the experts do not know how many more years the situation will stay the same. \\

% In Spain the recession has had a big impact that can still be seen in the unemployment numbers and jobs offers. Also, the poverty has increased enormously and the quality of life has decreased together with the salaries.  
% \\

% There is a general discontent with the politics and the banks. Pessimism is everywhere, as well as the feeling that it is needed a change of direction in both fields. 
% \\


% This economic context has created and motivated new ways of investigating and living, always aiming at spending as less as possible. It has triggered a change in the perception of the knowledge and whether it should be restricted using patents or not. Most of the scientific community now supports the open source initiative, a fact that has increased the rapidity of the new discoveries and technology improvements. 

% \begin{figure}[h]
% 	\begin{center}
%     \includegraphics[scale=0.3]{img/intro/osi.eps}
% 	\caption[Open Source Initiative Logo]{Open Source Initiative Logo}
% 	\end{center}
% \end{figure}


% The idea of creating common software, of aiding other investigators to easily replicate the work already done and continue from that point instead of "reinventing the wheel" started in the past century. In 1998 the Open Source Initiative (OSI)\cite{osi} was formed, and its definition is recognized as the standard. According to them, "Open source does not just mean access to the source code. The distribution terms of open-source software must comply with the following criteria: Free Redistribution, Source Code, Derived Works, Integrity of The Author's Source Code, No Discrimination Against Persons or Groups, No Discrimination Against Fields of Endeavor, Distribution of License, License Must Not Be Specific to a Product, License Must Not Restrict Other Software,  License Must Be Technology-Neutral"\cite{osi_def}. 
% \\

% Open Source is crucial in the development of new knowledge and not only in the Software field, but in all the science and technical disciplines. 



